\section{The Mesh}

The mesh is the whole space, one growing honeycomb, a hyperbolic tessellation in 4D. Its docks are stacked four-dimensional blocks, joined face to face, filling a curved space with no gaps. Its shape has a name, $\{3,4,3,4\}$, which the geometry part unpacks in full. For now three facts carry the weight. It is discrete, a countable weave of identical docks. It is hyperbolic, curved so that room opens up faster than it does in flat space. And it grows. Two words name its parts, and both recur throughout. The four-dimensional interior is the \emph{bulk}. The flat boundary it grows toward is the \emph{cusp}, the ordinary three-dimensional space we live in.

Why a honeycomb at all? Because the base must be discrete, and a regular honeycomb is the cleanest discrete space there is, every dock alike, every joint alike, no special places and no edges to fall off. Why hyperbolic, and not flat? Because feeling needs room to branch. In flat space a region grows only as fast as its surface. In hyperbolic space the interior swells, tree-like, with far more room a few steps out than any flat lattice allows. That extra room is what later lets a self hold a whole world inside itself, and what makes the boundary of a region behave like a screen. Which honeycomb, of the few that close up at all, is the subject of a later argument. The answer turns out to be forced, not chosen.

At any beat the mesh is finite, a definite patch of docks. It is never finished. New docks appear at its open edge each beat, so the space is unbounded over time though bounded at every instant. That open edge is where the wake works, and the steady growth there is the seed of time's direction. We come back to it in the law.
